\section{From propellant to a planning problem}
\label{sec:model}
The planning task is to move a spacecraft between prescribed orbital
states, possibly through planetary encounters, with finite propellant and
constant available useful jet power. The central question is how much
effort a specified arrival date requires. Once that function is available,
the earliest affordable arrival follows from a time search. This order
of calculation also exposes the marginal cost of changing dates and the
conditions needed to join independently represented trajectory phases.

All motion here is planar. Planet ephemerides are prescribed circular,
coplanar orbits; stations may orbit a planet or moon. Within each
gravitational patch the relative dynamics are controlled Kepler motion.
An unpowered point flyby and a finite planetary segment are alternative
encounter models. The ideal engine supplies useful analytic benchmarks;
the operational model uses full power or coast with exhaust speed in
the specified $50$--$250\,\mathrm{km/s}$ band. Observations in the
pursuit problem are instantaneous.

\subsection{The exact mass identity}
Let $P>0$ be constant useful jet power during a burn, $m$ the spacecraft
mass, $v_e$ exhaust speed and $a$ thrust acceleration magnitude. Then
\begin{equation}
 P=\tfrac12(-\dot m)v_e^2,\qquad ma=(-\dot m)v_e.
\end{equation}
Eliminating $v_e$ gives
\begin{equation}
 \begin{gathered}
 \frac{d}{dt}\frac1m=\frac{a^2}{2P},\qquad j=\int_0^T a^2dt,\\
 m_f=\left(\frac1{m_0}+\frac{j}{2P}\right)^{-1}.
 \end{gathered}
 \label{eq:mass}
\end{equation}
Coasting has $a=0$ and constant mass. Thus, for fixed initial mass and
power, minimizing propellant is exactly minimizing $j$. With dry mass
$m_d$, the available physical effort is
$B=2P(m_d^{-1}-m_0^{-1})$.
Electrical input power must first be converted to useful jet power.
The assumed $P$ does not increase as a fixed spacecraft becomes lighter.

For a primary with parameter $\mu$, choose a length $R_0$ and define
\begin{equation}
 \begin{gathered}
 T_0=\sqrt{R_0^3/\mu},\quad V_0=\sqrt{\mu/R_0},\\
 A_0=\mu/R_0^2,\quad J_0=V_0^2/T_0.
 \end{gathered}
\end{equation}
Dimensionless variables then obey
\begin{equation}
 \begin{gathered}
 r'=v,\qquad v'=g(r)+u,\\
 g(r)=-r/|r|^3,\qquad J=\int_0^h|u|^2dt,
 \end{gathered}
 \label{eq:kepler}
\end{equation}
where $h=T/T_0$ and $j=J_0J$. Define
\begin{equation}
 \begin{gathered}
 \chi=\frac{2P}{m_0J_0},\quad
 \chi_d=\frac{2P}{m_dJ_0},\quad
 \kappa=\chi_d-\chi,\\
 w=\frac{2P}{mJ_0}=\chi_d-b,\qquad
 w'=|u|^2,\quad b'=-|u|^2.
 \end{gathered}
 \label{eq:resource}
\end{equation}
Initially $b=\kappa$. In particular
\begin{equation}
 \frac{m_f}{m_0}=\frac{\chi}{\chi+J},\qquad
 J\le\kappa.
 \label{eq:massratio}
\end{equation}
This inverse-mass state makes resource propagation additive, including
across coasts and unpowered encounters.

\subsection{Ideal and bounded engines}
The ideal problem permits any square-integrable planar $u$. The
full-power/coast engine instead admits
\begin{equation}
 \begin{gathered}
 \mathcal U(b)=\{0\}\cup\{u:w/\nu_+\le|u|\le w/\nu_-\},\\
 \nu_\pm=V_{e,\pm}/V_0.
 \end{gathered}
 \label{eq:controlset}
\end{equation}
At $b=0$ only coast remains. A powered acceleration approaching zero
requires unbounded exhaust speed, so an ideal solution need not obey a
finite ceiling even when most of its burn does.

The ideal circular-endpoint action depends on $(\rho,\Theta,h)$, where
the departure radius is the length scale and $\rho$ is the arrival
radius ratio. Budget $\kappa$ determines affordability afterward.
The bounded action additionally depends on $\chi,\nu_-,\nu_+$ and the
available reserve. Fixing the band ratio removes one coordinate but
does not remove $\chi$. If $\mu,V_{e,+}$ and initial specific useful power
$\beta_{\rm jet}=P/m_0$ are all fixed across a vehicle family, the restricted
relation
\begin{equation}
 \chi=\frac{2\beta_{\rm jet}\mu}{V_{e,+}^{5}}\nu_+^5
\end{equation}
removes one more coordinate. Changing fuel load at fixed hardware power
does not preserve that family assumption.

\subsection{A free-space benchmark}
For $g=0$, fixed initial and final states and duration $h$, set
$d_r=r_f-r_0-hv_0$, $d_v=v_f-v_0$. The minimum effort is
\begin{equation}
 J_{\rm fs}=
 \frac{12|d_r|^2}{h^3}
 -\frac{12d_r\cdot d_v}{h^2}
 +\frac{4|d_v|^2}{h}.
 \label{eq:free}
\end{equation}
To see this, write the endpoint constraints as the two moments of $u$,
whose Gramian is
\begin{equation}
 W_0=\begin{pmatrix}h^3I/3&h^2I/2\\h^2I/2&hI\end{pmatrix}.
\end{equation}
The minimum-norm control is affine in time, and its squared norm is
$(d_r,d_v)^TW_0^{-1}(d_r,d_v)$, giving \eqref{eq:free}.
For a dimensional rest-to-rest displacement $L$,
\begin{equation}
 a(t)=\frac{6L}{T^2}\left(1-\frac{2t}{T}\right),\qquad
 j=\frac{12L^2}{T^3}.
\end{equation}
Consequently
$T_{\min}=[6L^2/(P(m_d^{-1}-m_0^{-1}))]^{1/3}$
in this ideal free-space problem.
Gravity changes the endpoint moments, so \eqref{eq:free} is not in
general a lower bound on Kepler effort.

\subsection{Provenance, scope and conventions}
The starting points are the author's constant-power Kepler transfer and
rocket pursuit--evasion articles \cite{OriginalKepler,OriginalChase}.
Inherited results include inverse-mass accounting, ideal free-space
transfers, the Kepler similarity reduction, interception and deadline-fuel
games, and the inertial gravity-compensation certificate. This article
develops the moving-endpoint derivatives $F_h,F_d$, the nonconvexity
example, the bounded resource adjoint on mass-dependent exhaust limits, the radial-measure
specialization, flyby and SOI junctions, the circular-spiral model, and the
CW rendezvous game. These are extensions or specializations of classical
optimal-control and game theory, not replacements for it
\cite{Pontryagin1962,Lawden1963,Isaacs1965}.
The planar restriction is deliberate. The chase article \cite{OriginalChase} also
treats $d=3$; the inertial free-space results below hold in either dimension.

Numerical data taken from the earlier articles come from repository commit
\nolinkurl{bdc9698212c92cd0815d9251c7d777d47cefd706}. All proofs needed for
the stated analytical results are included here. Numerical scripts and input
data are provided as supplementary material; a numerical observation is
distinguished from a proof.

\begin{center}\small
\begin{tabular}{@{}p{.32\columnwidth}p{.61\columnwidth}@{}}
\toprule Convention&Translation\\\midrule
Kepler time Hamiltonian&$H_T=\lambda\cdot f-1$, maximized; mass is a state\\
Ideal thrust gain&$C_m=\lambda_m m^2>0$, $k_T=P/C_m$\\
Physical effort costates&$p_{r,\rm phys}=2k_T\lambda_r$, $p_{v,\rm phys}=2k_T\lambda_v$\\
Normalized costates&$p_r=(R_0/J_0)p_{r,\rm phys}$, $p_v=(V_0/J_0)p_{v,\rm phys}$\\
Chase resource&$b_{\rm chase}=J_0b$ in $\mathrm{m^2/s^3}$\\
Kepler capability&$\kappa=\chi_d-\chi$\\
Mass ratio&$\eta_m=\chi_d/\chi=m_0/m_d$ (Kepler's $\eta$)\\
Exhaust limits&$\nu_\pm=V_{e,\pm}/V_0$\\
Minimizing DP convention&$p=-\nabla_s U$ along a differentiable optimal cost-to-go; also convert objective units\\
\bottomrule
\end{tabular}\end{center}
The costate gain conversion is for an ideal interior arc: there $k_T$ is
constant and $a=k_T\lambda_v=p_{v,\rm phys}/2$. Applying the same constant
conversion across a capped boundary is invalid. Multipliers for mass or
resource endpoints must be transformed with their state coordinate.

\begin{center}\small
\begin{tabular}{@{}p{.32\columnwidth}p{.61\columnwidth}@{}}
\toprule Symbol&Meaning and separation\\\midrule
$\beta_{\rm jet}$&$P/m_0$, W/kg; distinct from hardware coefficient $\beta_{\rm sys}$, kg/W\\
$Dg$&Jacobian of gravity; $D$ in the variational matrix is $BB^T/2$\\
$D(s)$&Flyby turn bound; $D$ in the interception section is separation beyond capture radius\\
$D$ in the spiral&Progress $|q_i-q_f|$ in Section~\ref{sec:stations}; distinct from the turn bound and matrix block\\
$C$&Resource coefficient $p_b+\ell_2$; $C(h)$ is a rendezvous effort; $\mathsf C$ denotes coast\\
$\omega$&Target mean motion in timing; local comparison frequency in verification, defined there; primer slope below is $n$\\
$L,\mathcal L$&Physical length (or a labeled lower bound), and the angular costate integral, respectively\\
$\eta_m,\eta$&Mass ratio and residual-error coefficient in the global-gap test, respectively\\
\bottomrule
\end{tabular}\end{center}


\subsection{Power--mass closure}
For design studies, let $m_d=m_{\rm pay}+\beta_{\rm sys}P$ and fix
$\mathcal R=m_{\rm prop}/m_d$. A useful hardware model is
\[
\beta_{\rm sys}=\frac{1}{(1-\varphi)\alpha_{\rm eng}}
+\frac{\varphi}{(1-\varphi)\rho_{\rm rad}},
\]
where $\varphi$ is the rejected fraction of input power,
$\alpha_{\rm eng}$ the engine specific input power and $\rho_{\rm rad}$
the radiator specific rejected power. Both specific powers are in W/kg.
Substituting $m_0=(1+\mathcal R)m_d$ into the ideal rest-to-rest formula gives
\begin{equation}
 T^3=\frac{6L^2(1+\mathcal R)}{\mathcal R}
          \left(\frac{m_{\rm pay}}P+\beta_{\rm sys}\right).
\end{equation}
Thus $T_\infty=[6L^2(1+\mathcal R)\beta_{\rm sys}/\mathcal R]^{1/3}$,
$P_{\rm knee}=m_{\rm pay}/\beta_{\rm sys}$, and
$T=T_\infty(1+P_{\rm knee}/P)^{1/3}$. This inherited hardware closure
\cite{OriginalKepler} is an ideal free-space design benchmark, not a
finite-band orbital prediction. It complements the fixed-hardware mass
identity without identifying W/kg with kg/W.
